%%%%%%%%%%%%%%%%%%%%%%%%%%%%%%%%%%%%%%%
% Cover Letter - Arcitecta, Student Internship (Oct 2026)
% Fun voice at Manit's request; prose only (DeakinTALENT rule).
%%%%%%%%%%%%%%%%%%%%%%%%%%%%%%%%%%%%%%%

\documentclass[]{cover}
\usepackage{fancyhdr}

\pagestyle{fancy}
\fancyhf{}

\rfoot{Page \thepage \hspace{0pt}}
\thispagestyle{empty}
\renewcommand{\headrulewidth}{0pt}
\begin{document}

%%%%%%%%%%%%%%%%%%%%%%%%%%%%%%%%%%%%%%
%     TITLE NAME
%%%%%%%%%%%%%%%%%%%%%%%%%%%%%%%%%%%%%%
\namesection{}{\Huge{Manit Khera}}{  \href{mailto:manitkhera26@gmail.com}{manitkhera26@gmail.com} | 0412 761 206 |  \urlstyle{same}\href{https://www.linkedin.com/in/manit-khera-603957252}{LinkedIn}
}

%%%%%%%%%%%%%%%%%%%%%%%%%%%%%%%%%%%%%%
%     MAIN COVER LETTER CONTENT
%%%%%%%%%%%%%%%%%%%%%%%%%%%%%%%%%%%%%%

\currentdate{20 October 2026}
\lettercontent{Dear Hiring Manager,}

\lettercontent{I found your Student Internships on the Arcitecta careers page, and one line felt written for me: you want makers who will step into problems “even when there are many unknowns”. I go looking for those on purpose.}

\lettercontent{I am a third-year student in a Bachelor of Software Engineering (Honours) at Deakin University, finishing in 2028, with a WAM of 81.33 and 100 in Data Structures and Algorithms. My backend unit chair had a habit of telling me “nah, you can't do this”, which works on me every time. So I made my C\# REST API drive a robot I had wired myself. It scored 100/100 and I demonstrated it live at Deakin Open Day, which fits your “regularly demonstrate working software” line.}

\lettercontent{You also ask for pragmatism over dogma. My robot only behaved on fresh batteries. Rather than guess, I read its own boot and reset logs, which blamed a failing buck converter, not my code. By then I had bought a small pile of batteries that are now all about 90 percent full. On concurrency, I audited a robot dispatch API for race conditions and fixed them with atomic PostgreSQL job claims, so each job is claimed exactly once.}

\lettercontent{Mediaflux is the kind of software I want to learn from. The National Film and Sound Archive chose it as its central digital asset management system, and Mediaflux Real-Time won a 2025 NAB Show Product of the Year award. You own your entire technology architecture, and I like working across layers too: my robot runs from Arduino firmware up to a C\# REST API.}

\lettercontent{One gap: I have not written Java yet. I work mostly in C\#, plus C++ for Arduino firmware and JavaScript for React, and you say you will help, so I will hold you to that. I am on a student visa, so I can work full-time from November to February and two to three days a week during trimester.}

\lettercontent{My resume is attached. I would love an interview, ideally with the robot in tow. Thanks for reading!}

\begin{flushright}
\closing{Kind regards,}

\signature{Manit Khera}
\end{flushright}
\end{document}
